\textit{Note: This section is out of syllabus for NSW HSC Mathematics Extension~1 and Extension~2. It is included for enrichment and to give a first glimpse of higher-dimensional calculus.}

\subsubsection*{Partial differentiation}
In ordinary differentiation, a function depends on one independent variable, such as $y=f(x)$. In higher mathematics, functions often depend on several variables at once.

For example, a temperature function might be written as
\[
T=T(x,y,z,t),
\]
where temperature depends on position and time.

When differentiating with respect to one variable while treating the others as constants, we use a \textbf{partial derivative}. The notation is
\[
\frac{\partial f}{\partial x}.
\]

If
\[
f(x,y)=x^3y^2+5x,
\]
then
\[
\frac{\partial f}{\partial x}=3x^2y^2+5,
\qquad
\frac{\partial f}{\partial y}=2x^3y.
\]

\subsubsection*{Partial differential equations}
A \textbf{partial differential equation} (PDE) relates a multivariable function to one or more of its partial derivatives.

For example, the one-dimensional heat equation is
\[
\frac{\partial u}{\partial t}=\alpha \frac{\partial^2u}{\partial x^2},
\]
where $u(x,t)$ is temperature and $\alpha$ is a constant.

This equation says that the rate of change of temperature over time depends on how curved the temperature profile is in space.

\subsubsection*{Why PDEs matter}
PDEs are central in heat flow, fluid motion, wave propagation, electromagnetism, and quantum mechanics. They sit well beyond the HSC course, but they grow naturally from the same calculus ideas students already know.
